\section{Reconstruction and the actual HYM flow}\label{sec:actual-flow}

The base and bundle parameters are fixed throughout
this section. All exact initial metrics are smooth,
have the parallel determinant, and obey fixed positive
bounds against $h_0$. In local unitary flat frames,
these conditions mean $\det H=1$ and $cI\le H\le CI$.
We use physical time $s$ and the normalization
\begin{equation}\label{eq:hym}
 H^{-1}\partial_sH=-2i\Lambda_{\omega_\tau}F_H.
\end{equation}
For $\omega_\tau=(i/2)a\,dz\wedge d\bar z$ this reads
\begin{equation}\label{eq:local-hym}
 \partial_sH=4a^{-1}
     (H_{z\bar z}-H_{\bar z}H^{-1}H_z).
\end{equation}
The scalar nonnegative Laplacian is
$\Delta_\tau=-4a^{-1}\partial_z\partial_{\bar z}$.
The real-coordinate expression is
\begin{equation}\label{eq:real-hym}
 \begin{split}
 a\partial_sH={}&H_{xx}+H_{yy}-H_xH^{-1}H_x-H_yH^{-1}H_y\\
 &+i(H_xH^{-1}H_y-H_yH^{-1}H_x).
 \end{split}
\end{equation}
The last term need not vanish. We retain it in every
curvature calculation.

For two metrics put
\[
 \sigma(H,K)=\tr(H^{-1}K+K^{-1}H)-2r.
\]
This is intrinsic under the bundle transition maps.
If $e^{b_j}$ are the eigenvalues of $H^{-1}K$, then
$d(H,K)^2=\sum_jb_j^2\le
\sum_j(e^{b_j}+e^{-b_j}-2)=\sigma(H,K)$.
For a smooth comparison metric define
\begin{equation}\label{eq:residual}
 \mathcal R(K)=K^{-1}\partial_sK+2i\Lambda F_K,
\end{equation}
measured in the operator norm induced by $K$.

\begin{lemma}\label{lem:trace-comparison}
The exact solution of \eqref{eq:hym} exists for all
$s\ge0$ and preserves its determinant and bounds
$ch_0\le H\le Ch_0$. For any smooth comparison
$K(s)$ with common positive bounds, one has
\begin{equation}\label{eq:trace-comparison}
 \int\sigma(H(s),K(s))
 \le\int\sigma(H(0),K(0))
       +C_0\int_0^s\!\int|\mathcal R(K(v))|_{\op}\,dA_\tau\,dv .
\end{equation}
Here and below unlabelled spatial integrals use $dA_\tau$.
The constant depends only on rank and the positive
bounds, not on a spectral gap or Sobolev constant.
\end{lemma}
\begin{proof}
Smooth all-time existence on a fixed compact base is
the classical metric heat-flow result
\cite{donaldson}; our normalization and order
preservation are also given in
\cite[Section~4.1 and Proposition~4.3]{hkkp}.
Apply order preservation to the stationary metrics
$ch_0,Ch_0$. The logarithm of the determinant satisfies
the scalar heat equation, so its constant initial
value is preserved.

For completeness the local Donaldson trace calculation
has the forced form, with $q=H^{-1}K$,
\[
 (\partial_s+\Delta_\tau)\sigma(H,K)
 \le\tr((q-q^{-1})\mathcal R(K)).
\]
For two exact velocities choose a holomorphic frame
normal for $H$ at the point. The curvature and
Laplacian terms in the forward trace cancel, leaving
$-4a^{-1}\tr(K_{\bar z}K^{-1}K_z)\le0$.
Interchange the metrics for the inverse trace.
The additional velocity $K\mathcal R(K)$ contributes
$\tr(q\mathcal R)$ and $-\tr(\mathcal Rq^{-1})$,
respectively. Both $q$ and $\mathcal R$ are
$K$-self-adjoint, and the common bounds control
$\tr q+\tr q^{-1}$. This proves the scalar inequality
with right-hand side at most $C_0|\mathcal R|_{\op}$.
Integration on the whole closed fibre removes the
Laplacian and gives \eqref{eq:trace-comparison}.
The calculation is unchanged by constant unitary
transition maps.
\end{proof}

Write $\varepsilon_e=L_e^{-1}$,
$\varepsilon_*=L_*^{-1}$, and
\begin{equation}\label{eq:error-scale}
 e_*=\frac{(1+\log L_*)^2}{L_*}.
\end{equation}
The initial tuples for \eqref{eq:finite-flow} lie in
a fixed bounded ball.

\begin{theorem}\label{thm:tracking}
For all sufficiently large $L_*$ there is a smooth
determinant-one comparison metric $\widehat K(s)$,
defined for all $s\ge0$ from the finite trajectory
\eqref{eq:finite-flow}, with uniform positive bounds
and
\begin{equation}\label{eq:residual-bound}
 \sup_{s\ge0}\int|\mathcal R(\widehat K(s))|_{\op}
 \le C L_*^{-2}(1+\log L_*)^2.
\end{equation}
On every fixed compact subset of component $i$,
$d(\widehat K(s),H_i(s))\le C_K/L_*$ uniformly in
$s$. If an exact flow has initial discrepancy
$\eta=\int\sigma(H^{\mathrm{in}},\widehat K(0))$,
then, for every $T\ge0$,
\begin{equation}\label{eq:tracking}
 \sup_{0\le s\le TL_*}\int
 d(H(s),\widehat K(s))^2\le\eta+CTe_*.
\end{equation}
The constants are independent of $T$, all positive
length ratios and all unitary parameters. Uniform
convergence over a class of initial metrics requires
uniform $\eta\to0$.
\end{theorem}

\subsection{Component potentials}
Let $g_{ie}$ be the mean-zero scalar Green function
on the fixed closed component $C_i$, normalized by
\[
 \Delta_i g_{ie}=-2\pi\delta_{p_{ie}}+\frac{2\pi}{A_i},
 \qquad g_{ie}=\log|u_e|+\gamma_{ie}(u_e)
\]
near the incidence. The regular part $\gamma_{ie}$
is smooth. Define
\begin{equation}\label{eq:component-correction}
 h_i=-\sum_{e=ij}\varepsilon_e B_{ij}g_{ie},
 \qquad K_i=H_i\exp(h_i).
\end{equation}
The logarithm $h_i$ is trace free and $H_i$-self-adjoint.
Away from the punctures,
\begin{equation}\label{eq:cancellation}
 \Delta_i h_i=-H_i^{-1}\dot H_i.
\end{equation}
Cut the components at radius $\rho=\varepsilon_*$.
Since $\max|\tau_e|=e^{-L_*}\ll\rho^2$, the two
branches of each node are disjoint there.
Fixed Green-function estimates and
Proposition~\ref{prop:finite-flow} give
\begin{equation}\label{eq:green-bounds}
 \begin{aligned}
 \|h_i\|_\infty&\le C\varepsilon_*(1+|\log\rho|),
 &\int|h_i|&\le C\varepsilon_*,
 &\int|h_i|^2&\le C\varepsilon_*^2,\\
 \int|\nabla h_i|^2&\le C\varepsilon_*^2(1+|\log\rho|),
 &\int|\dot h_i|&\le C\varepsilon_*^2,
 &\int|h_i||\Delta_i h_i|&\le C\varepsilon_*^2 .
 \end{aligned}
\end{equation}
In particular the integral bounds use integrability
of $\log|u|$ and its square, not their suprema on
the truncated domain. The time derivative estimate
uses $|\dot B_{ij}|\le C\varepsilon_*$ and
$\varepsilon_e\le\varepsilon_*$.

The exact exponential identity
\[
 e^{-h}\partial e^h
 =\int_0^1e^{-th}(\partial h)e^{th}\,dt
\]
and its time analogue show, after
\eqref{eq:cancellation}, that
\[
 |\mathcal R(K_i)|
 \le C\big(
 |\dot h_i|+\varepsilon_*|h_i|
 +|\nabla h_i|^2+|h_i||\Delta_i h_i|\big).
\]
Indeed the linear spatial term is $\Delta_i h_i$;
all terms produced by differentiating the conjugations
contain either two first derivatives or $h_i$ times
a second derivative. The time conjugation of
$H_i^{-1}\dot H_i$ contributes
$O(\varepsilon_*|h_i|)$. These identities include
the commutators in \eqref{eq:real-hym}. Integrating
\eqref{eq:green-bounds} bounds the residual by
$C\varepsilon_*^2(1+|\log\rho|)$ on the fixed
component. Transfer to the actual exterior costs
$Cb\rho^{-N}$, where $b=e^{-L_*}$, by
Lemma~\ref{lem:ambient}. A fixed $N$ suffices for
the finite derivative order involved.

\subsection{Node profiles and mixed weights}
Near an incidence write
\[
 h_i=-\varepsilon_eB_{ij}\log|u_e|+c_{ie}+q_{ie}(u_e),
 \qquad q_{ie}(0)=0.
\]
The regular terms include contributions from every
incident edge. Uniformly on the collar,
$|c_{ie}|\le C\varepsilon_*$ and
$|D^kq_{ie}|\le C\varepsilon_*\rho^{1-k}$
for $k\le2$. Their time derivatives gain
a factor $\varepsilon_*$.
Set $P_{ie}=H_i\exp(c_{ie})$. In the $i$ frame
the node metric is
\begin{equation}\label{eq:node-profile}
 N_e(s,u_e)=
 \Geo(P_{ie},U_{ij}P_{je}U_{ij}^*;
                  -\varepsilon_e\log|u_e|).
\end{equation}
The reverse branch has parameter
$1+\varepsilon_e\log|u_e|$, because $u_ev_e=\tau_e$.
Geodesic reversal and equivariance under unitary
congruence therefore give exactly the same metric
in the other frame. The individual $\varepsilon_e$
is essential for this compatibility.

At fixed time a geodesic in
\eqref{eq:node-profile} is $P\exp(\theta B)$
with a fixed matrix $B$ and harmonic scalar
$\theta$. Its Chern curvature is zero on the
annulus. Its time derivative is $O(\varepsilon_*)$
on the bounded endpoint range; the middle region
has area $O(\rho^2)$, so its integrated time
residual is $O(\varepsilon_*\rho^2)$.

On $\rho\le|u_e|\le2\rho$ use relative logarithms
$\ell_i(K)=\log(H_i^{-1}K)$.
For fixed endpoint metrics $H_i,H_j$, and
small self-adjoint endpoint logarithms $C_i,C_j$,
define the difference
\[
 \mathcal D(C_i,C_j;\theta)
 =\ell_i\!\left(\Geo(H_i e^{C_i},
              U_{ij}H_j e^{C_j}U_{ij}^*;\theta)\right)
                 -C_i-\theta B_{ij}.
\]
Self-adjointness here is with respect to
$H_i$ and $H_j$, respectively.
With no endpoint correction the node logarithm
is exactly $\theta B_{ij}$, so
$\mathcal D(0,0;\theta)=0$. At $\theta=0$
the geodesic equals $H_i e^{C_i}$, so
$\mathcal D(C_i,C_j;0)=0$.
The geodesic and logarithm are smooth on the
confined endpoint range. Integrating their
bounded mixed derivatives first along
$q(C_i,C_j)$, $0\le q\le1$, and then from
$0$ to $\theta$ gives
\[
 |\mathcal D(C_i,C_j;\theta)|
       \le C(|C_i|+|C_j|)|\theta|.
\]
Its first two $\theta$ derivatives are bounded
by $C(|C_i|+|C_j|)$, since they also vanish
when both corrections are zero.

Apply this with $C_i=c_{ie}$, $C_j=c_{je}$
and $\theta=-\varepsilon_e\log|u_e|$.
The endpoint corrections are $O(\varepsilon_*)$,
whereas the spatial derivatives of $\theta$
have size $O(\varepsilon_e\rho^{-k})$
for $k=1,2$. Thus the correction to the node
logarithm is bounded, through two spatial
derivatives, by
$C\varepsilon_*\varepsilon_e
(1+|\log\rho|)\rho^{-k}$.
The component logarithm differs from
$c_{ie}+\theta B_{ij}$ by the regular term
$q_{ie}$, whose corresponding bound is
$C\varepsilon_*\rho^{1-k}$.
Combining the two estimates gives
\begin{equation}\label{eq:mixed-collar}
 |D^k(\ell_i(K_i)-\ell_i(N_e))|
 \le C\left(
 \varepsilon_*\rho^{1-k}
 +\varepsilon_*\varepsilon_e(1+|\log\rho|)\rho^{-k}
 \right),\qquad k=0,1,2.
\end{equation}
The endpoint correction uses all incident weights,
whereas the node slope uses $\varepsilon_e$.
No inverse of $\varepsilon_e$ enters.

Join these logarithms with a fixed rescaled radial
cutoff and exponentiate relative to $H_i$. This
preserves positivity and determinant one. If $n$
is the node logarithm and $\delta$ is its cutoff
difference from the component logarithm, the full
curvature difference is bounded by
\[
 C\big(
 |D^2\delta|+|Dn||D\delta|
 +|\delta|(|D^2n|+|Dn|^2)+|D\delta|^2
 \big).
\]
Here $|Dn|\le C\varepsilon_e/\rho$ and
$|D^2n|\le C\varepsilon_e/\rho^2$.
Multiply by the collar area $O(\rho^2)$ and
substitute \eqref{eq:mixed-collar}. The integrated
curvature error is at most
\[
 C\big(\varepsilon_*\rho+
       \varepsilon_*\varepsilon_e(1+|\log\rho|)^2\big).
\]
Time terms contribute $O(\varepsilon_*\rho^2)$.
Terms of higher order are absorbed since
$\varepsilon_*|\log\rho|\to0$. The same geometric
transfer estimate contributes $Cb\rho^{-N}$.

\begin{proof}[Proof of Theorem~\ref{thm:tracking}]
Use \eqref{eq:component-correction},
\eqref{eq:node-profile} and the logarithmic cutoffs
to define one global $\widehat K(s)$ for the
entire finite trajectory. Their agreement under
the actual transition matrices proves that this
is a metric on $E_U$, not just a collection of
matrices. The confined endpoint range and the
small component logarithms give uniform positive
bounds. On each fixed component compact subset
$h_i=O(\varepsilon_*)$.

Sum the component, middle-node and collar estimates
with $\rho=\varepsilon_*$. The exponentially small
$e^{-L_*}L_*^N$ is absorbed by
$L_*^{-2}(1+\log L_*)^2$. All constants are
uniform for $s\ge0$, because the finite trajectory
stays in one compact range. This proves
\eqref{eq:residual-bound}.
Lemma~\ref{lem:trace-comparison} on the whole
closed fibre then gives \eqref{eq:tracking}
after integration to $TL_*$ and the inequality
$d^2\le\sigma$.
\end{proof}

\subsection{A uniform all-time conclusion}
\begin{theorem}\label{thm:all-time}
Under the hypotheses of Theorem~\ref{thm:tracking},
suppose in addition that
\begin{equation}\label{eq:rate-condition}
 L_*\nu_0\ge\kappa>0.
\end{equation}
Then, for sufficiently large $L_*$,
\[
 \sup_{s\ge0}\int d(H(s),\widehat K(s))^2
 \le C\eta+C_\kappa e_*(1+|\log e_*|).
\]
This is absolute tracking in the indicated norm.
\end{theorem}
\begin{proof}
For the exact solution set
$\mathcal F(H)=\int(\tr(h_0^{-1}H)-r)$.
Since $S=\log(h_0^{-1}H)$ is trace free and the
metrics have fixed positive bounds,
$\mathcal F$ is uniformly comparable to
$\int|S|^2=\int d(H,h_0)^2$. Taking the trace
of \eqref{eq:local-hym} and integrating gives
\[
 -\mathcal F'
 =4\int\tr(H_{\bar z}H^{-1}H_z)\,dx\,dy
 \ge c_1\int|D_UH|^2\,dA_\tau .
\]
Here $H_{\bar z}=H_z^*$. Replace $H^{-1}$ by
its uniform lower bound before using
$4\tr(H_z^*H_z)=\tr(H_x^2+H_y^2)$.
This does not discard a weighted commutator.
The derivative of the matrix logarithm is bounded
on the common matrix range, and hence
$|D_US|\le C|D_UH|$. Theorem~\ref{thm:spectrum}
and \eqref{eq:rate-condition} imply
$\lambda_0\ge\kappa/(2L_*)$ and therefore
\[
 -\mathcal F'\ge c_2\lambda_0\mathcal F,\qquad
 \int d(H(L_*\sigma),h_0)^2
       \le C e^{-c_3\kappa\sigma}.
\]

Lemma~\ref{lem:finite-coercivity} supplies the
corresponding exponential tail for the finite
trajectory. The reconstruction vanishes at
$\mathbf H=\mathbf I$ and is uniformly Lipschitz
from the confined tuple range to global $L^2$:
on components this follows from the Green
construction and on nodes from smoothness of
bounded geodesics. Its squared distance from
$h_0$ has the same type of exponential tail.
For $0\le s\le TL_*$ use \eqref{eq:tracking};
for $s\ge TL_*$ use the two tails and the
triangle inequality. Choosing
$T=C_\kappa(1+|\log e_*|)$ proves the estimate.
\end{proof}

The condition \eqref{eq:rate-condition} permits
some arbitrarily weak edges. In rank two keep
$01,02,03,12,13$ comparable to $w_*$, with the
two noncommuting transports of
Example~\ref{ex:noncommuting}, and let $23$
have any smaller positive weight. The first
five edges already have zero trace-free
kernel and a positive normalized gap; the
last nonnegative term cannot decrease it.
In contrast, Proposition~\ref{prop:three-scales}
does not satisfy \eqref{eq:rate-condition}.
On its weakest time $h^{-4}$ the accumulated
bound in \eqref{eq:tracking} is not small.
Uniform nonlinear tracking on that time
requires an additional argument.
